\KOMAoptions{parskip=false}

%% Margins -----------------------------
\usepackage[a4paper,
    head=\baselineskip,
    top=25.8mm,
    bottom=24.8mm,
    inner=40mm,
    outer=21mm
]{geometry}

% Footers -----------------------------
\setlength{\footheight}{11.7mm}
\clearscrheadfoot
\ifoot*{\Title\vfill}
\ofoot*{\pagemark\vfill}
\setkomafont{pageheadfoot}{\fontsize{9pt}{13pt}\normalfont}
\setkomafont{pagefoot}{\bfseries}
\setkomafont{pagenumber}{\normalfont}
\pagestyle{scrheadings}

% Footnotes ---------------------------
%  only use if really necessary - most likely, you want should use a reference instead
%\KOMAoptions{%
%    footnotes=multiple % mehrere Fußnoten werden durch Zeichen getrennt
%}
%\setfootnoterule[.6pt]{5.08cm}
%\renewcommand{\footnoterule}{\hrule width 5.08cm height .6pt \vspace*{3.9mm}}
%\setlength{\footnotesep}{2mm}
%\deffootnote{2mm}{2mm}{%
%    \makebox[2mm][l]{\textsuperscript{\thefootnotemark}}%
%}
%\setkomafont{footnoterule}{\fontsize{9pt}{20pt}\selectfont}


% Headings ----------------------------
\KOMAoptions{%
    open=any,
    numbers=noendperiod, % no punctation after numbers
    headings=small
}

\makeatletter
\g@addto@macro{\@afterheading}{\vspace{-\parskip}} % remove \parskip after headings
\renewcommand*{\chapterheadstartvskip}{\vspace{\@tempskipa}\vspace{-3pt}} % correct space between headlines
\makeatother

\setkomafont{disposition}{\normalfont\sffamily}

\setkomafont{chapter}{\normalfont\fontsize{19pt}{22pt}\selectfont}
\RedeclareSectionCommand[%
  beforeskip=0pt,
  afterskip=29pt
]{chapter}
\renewcommand*{\chapterformat}{\thechapter.\enskip} % add punctation after chapter numbers

\setkomafont{section}{\fontsize{15pt}{17pt}\selectfont}
\RedeclareSectionCommand[%
  beforeskip=0pt,
  afterskip=24.1pt
]{section}
\renewcommand*{\sectionformat}{\makebox[13mm][l]{\thesection.\enskip}}

\setkomafont{subsection}{\bfseries\fontsize{12pt}{13pt}\selectfont}
\RedeclareSectionCommand[%
  beforeskip=0pt,
  afterskip=1pt
]{subsection}
\renewcommand*{\subsectionformat}{\makebox[13mm][l]{\thesubsection.\enskip}}

\setkomafont{subsubsection}{\itshape\fontsize{11pt}{12pt}\selectfont}
%\RedeclareSectionCommand[%
%  beforeskip=0pt,
%  afterskip=1pt
%]{subsubsection}
\renewcommand*{\subsubsectionformat}{\makebox[13mm][l]{\thesubsubsection.\enskip}}

% Lists
\setlist[description]{%
    labelindent=!,
    leftmargin=1em,
    labelwidth=!,
    parsep=0.5\parskip,
    partopsep=0mm,
    labelsep=1em,
}

%% List of ... -------------------------
\KOMAoptions{%
%    toc=flat, % keine Einrückungen im Inhaltsverzeichnis
    toc=chapterentrydotfill, % Punkte bis zur Seitennummer bei Kapiteln
    listof=entryprefix, % Präfix für Einträge in Abbildungs- und Tabellenverzeichnis
%    listof=nochaptergap, % Kein Abstand für Kapiteleinträge in extra Verzeichnissen
}

\makeatletter
\renewcommand{\@dotsep}{.3} % Abstand der Füllpunkte

% "chapteratlist" für Inhaltsverzeichnis auswerten:
\renewcommand*{\addchaptertocentry}[2]{%
  \iftocfeature{toc}{chapteratlist}{}{%
    \addtocontents{toc}{\protect\vspace{-5pt}}% extra Abstand vor Kapitelüberschriften in Inhaltsverzeichnis entfernen
  }%
  % Originaldefinition aus scrbook.cls:
  \addtocentrydefault{chapter}{#1}{#2}%
  \if@chaptertolists
    \doforeachtocfile{%
      \iftocfeature{\@currext}{chapteratlist}{%
        \addxcontentsline{\@currext}{chapteratlist}[{#1}]{#2}%
      }{}%
    }%
    \@ifundefined{float@addtolists}{}{\scr@float@addtolists@warning}%
  \fi%
}
\makeatother

\AfterTOCHead[toc]{\protect\vspace{.8ex}} % Abstand zwischen Überschrift und Inhaltsverzeichnis
\setuptoc{toc}{noparskipfake} % Angleichung der Abstände nach Inhaltsverzeichnisüberschrift an andere Überschriften
\unsettoc{toc}{chapteratlist} % kein Abstand vor Kapiteleinträgen im Inhaltsverzeichnis, funktioniert nur durch obige Redefinition von \addchaptertocentry

% Figures and Tables
\AfterTOCHead[lof]{\protect\vspace{-.1ex}\doublespacing} % Abstand zwischen Überschrift und Abbildungsverzeichnis, doppelter Zeilenabstand
\setuptoc{lof}{noparskipfake} % Angleichung der Abstände nach Abbildungsverzeichnisüberschrift an andere Überschriften

\AfterTOCHead[lot]{\protect\vspace{-.1ex}\doublespacing} % Abstand zwischen Überschrift und Tabellenverzeichnis, doppelter Zeilenabstand
\setuptoc{lot}{noparskipfake} % Angleichung der Abstände nach Tabellenverzeichnisüberschrift an andere Überschriften

% Caption
\DeclareCaptionFormat{WissenschaftlicheArbeiten}{\fontsize{10pt}{11pt}\selectfont#1 #2#3\par}
\DeclareCaptionLabelFormat{WissenschaftlicheArbeiten}{\bfseries\selectfont#1 #2}
\captionsetup[table]{%
    format=WissenschaftlicheArbeiten,
    labelformat=WissenschaftlicheArbeiten,
    labelsep=none,
    singlelinecheck=off,
    justification=centering,
    skip=3pt,
    tablewithin=none
}
\captionsetup[figure]{%
    format=WissenschaftlicheArbeiten,
    labelformat=WissenschaftlicheArbeiten,
    labelsep=none,
    singlelinecheck=off,
    justification=raggedright,
    skip=6.6mm,
    figurewithin=none
}

% Tables
\renewcommand{\arraystretch}{1.0} % Skalierung der Tabellen
\newcolumntype{M}{X<{\vspace{4pt}}} % Spaltentyp mit Abstand rechts
\newcolumntype{C}[1]{>{\centering\arraybackslash}m{#1}}